\documentclass[a4paper,11pt,fleqn]{article}%fleqn pour aligner les equations à gauche
\usepackage{../../Styles/Cours}

\newcommand{\chnum}{Chapitre 11}
\newcommand{\chtitle}{Caractéristiques d'un mouvement}
\newcommand{\dtype}{TP}



\begin{document}
\maketitle

Pour décrire un mouvement, on utilise deux adjectifs : l'un indique la forme de la trajectoire et l'autre informe sur l'évolution de la vitesse au cours du temps. Ainsi, le mouvement d'un système qui suit un cercle à vitesse constante est qualifié de \textit{circulaire} et \textit{uniforme}.\\ Il est possible d'étudier les caractéristiques du mouvement d'un système à partir d'une vidéo.
\section{Trajectoire}
	\subsection*{Position du système}
	\begin{enumerate}
		\item Ouvrir le fichier vidéo fourni avec le logiciel \textit{Latis Pro} puis lire la vidéo.
		\item Régler les paramètres d'étude :
		\begin{itemize*}
			\item sélection de l'origine : sélectionner le bas de la fenêtre vidéo
			\item sélection de l'étalon : sélectionner la règle et indiquer sa taille (\SI{0,507}{m})
			\item sens des axes : vers le haut
			\item déplacement absolu
			\item sélection manuelle des points
		\end{itemize*}
		\item Repérer les positions successives du système en cliquant dessus image après image
		\item Enregistrer le fichier
		\item Transférer vers les vecteurs. Un document s'ouvre représentant la trajectoire du système. Comment peut-on qualifier cette trajectoire ?
	\end{enumerate}
	\subsection*{Vitesse}
	\begin{enumerate}[resume]
	\item Rappeler les caractéristiques du vecteur vitesse $\vec{v_i}$
	\item Calculer les valeurs des vecteurs vitesse aux points $M_2$, $M_3$, $M_6$ et $M_7$
	\item Les représenter sur la trajectoire à l'aide d'une échelle de votre choix.
	\item Comment évolue le vecteur vitesse au cours du temps ?
	\end{enumerate}
	\subsection*{Accélération}
	L'accélération correspond à la variation de la vitesse par rapport au temps.
	\begin{enumerate}[resume]
	\item Donner l'expression du vecteur accélération $\vec{a_i}$
	\item Représenter les vecteurs variation de vitesse aux points $M_2$ et $M_6$
	\item Calculer les valeurs des vecteurs accélération aux points $M_2$ et $M_6$
	\item Les représenter sur la trajectoire à l'aide d'une échelle choisie
	\item Comment évolue le vecteur accélération au cours du temps ?
	\item Qualifier le mouvement
	\end{enumerate}
\section{Coordonnées}

La position d'un point $M$ à la date $t$ est donnée par le vecteur $\vv{OM}(t)$, de coordonnées :\\
	\begin{tabular}{>{\centering}m{6cm} | m{13cm}}
	\multirow{2}{5cm}{\begin{eqnarray}
	\vv{OM}(t) \left\{ 
	\begin{array}{l} x(t) \\ y(t) \\ z(t) \\ 
	\end{array} \right\} _{(O, \vec{i}, \vec{j}, \vec{k})} 	\nonumber
	\end{eqnarray}} & $x(t)$ : abscisse du point M (en $m$)\rule[.75cm]{0cm}{0cm}\\
	& $y(t)$ : ordonnée du point M (en $m$)\rule[.75cm]{0cm}{0cm}\\
	& $z(t)$ : altitude du point M (en $m$)\rule[.75cm]{0cm}{0cm}\\
	\end{tabular}\\
Le vecteur vitesse $\vec{v}(t)$ de $M$ à la date $t$ est la dérivée du vecteur position par rapport au temps : $\vec{v}=\dfrac{d \vv{OM}}{dt}$\\
Il a pour coordonnées (en \unit{\meter\per\second}) :\\
$$\vec{v}(t)
	\left( 
	\begin{array}{l} v_x(t)=\dfrac{dx}{dt} \\ v_y(t)=\dfrac{dy}{dt} \\ v_z(t)=\dfrac{dz}{dt} \\ 
	\end{array} 
	\right) _{(O, \vec{i}, \vec{j}, \vec{k})}$$
\begin{enumerate}[resume]
	\item Ouvrir le tableur
	\item Glisser la courbe liée à $x$ dans une colonne et celle liée à $y$ dans une autre. Comment évolue la position $x$ du système au cours du temps ? et sa position $y$ ?
	\item Représenter l'évolution de la coordonnée $x(t)$ en fonction du temps et l'évolution de la coordonnée $y(t)$ en fonction du temps
	\item Modéliser les courbes obtenues et relever leurs équations :\\
	$x(t)$ =\\
	$y(t)$ =
	\item À l'aide du document ci-dessous et de la fiche maths sur les dérivées, déterminer l'expression des coordonnées du vecteur vitesse :\\
	$v_x(t)$ =\\
	$v_y(t)$ =
	\item En déduire l'expression des coordonnées du vecteur accélération : \\
	$a_x(t)$ =\\
	$a_y(t)$ =
	\item Utiliser les expressions précédentes pour donner les coordonnées des vecteurs vitesse et accélération au point $M_2$ (date $t_2$).\\
	Comparer ces coordonnées et les représentations graphiques de la partie \textbf{1}.
\end{enumerate}
\end{document}
